% core/stone-onepage.tex — 罗塞塔石碑（一页版），独立编译
\documentclass[9pt]{article}
\usepackage[UTF8, fontset=windows]{ctex}
\usepackage[a4paper,landscape,margin=1.1cm]{geometry}
\usepackage{amsmath,amssymb,bm}
\usepackage{xcolor}
\usepackage{array,longtable,booktabs}
\usepackage{enumitem}
\definecolor{mcol}{RGB}{28,78,140}
\definecolor{pcol}{RGB}{150,74,10}
\definecolor{gcol}{RGB}{105,105,105}
\linespread{1.05}
\pagestyle{empty}
\setlength{\parindent}{0pt}

\begin{document}

{\Large\bfseries 罗塞塔石碑：同一件事的四种语言}\hfill
{\small\color{gcol}Kapustin--Witten 2006 · 几何 Langlands 纲领 · 数论}

\vspace{2pt}
{\small\color{gcol}
韦伊（A. Weil）说他的工作像在破译一块三语石碑：数论、$\mathbb F_q$ 上的曲线、黎曼面，三栏彼此平行；
他这些年只找到「词典的一小块一小块」。KW06 加上第四栏——四维 $\mathcal N=4$ 规范理论。
下表每行是同一件事在四栏里的名字。}

\vspace{6pt}
\begingroup\footnotesize
\setlength{\tabcolsep}{4pt}
\renewcommand{\arraystretch}{1.28}
\begin{longtable}{@{}p{2.5cm}p{4.5cm}p{5.3cm}p{5.3cm}p{4.6cm}@{}}
\toprule
\textbf{核心组件} & \textbf{数论} & \textbf{$\mathbb F_q$ 上的曲线} & \textbf{黎曼面（几何侧）} & \textbf{4d $\mathcal N=4$（物理侧）} \\
\midrule
\endhead
基域 / 舞台
& $\mathbb Q$ 与其有限扩张
& 函数域 $\mathbb F_q(C)$
& 紧黎曼面 $C$（亏格 $g$）
& 紧化曲线 $C$ 上的 4d 理论 \\
\midrule
对偶群
& Langlands 对偶群 ${}^{L}G$
& 同左
& 平坦 ${}^{L}G$-丛（局部系统）
& $S$-对偶后的规范群 ${}^{L}G$ \\
\midrule
耦合 / 模参数
& —（无耦合）
& —（无耦合）
& $\tau$：$C$ 的复结构模（Teichmüller）
& $\tau=\dfrac{\theta}{2\pi}+\dfrac{4\pi i}{g^{2}}\in\mathbb H$ \\
\midrule
谱侧对象
& Galois 表示 $\mathrm{Gal}(\bar{\mathbb Q}/\mathbb Q)\to{}^{L}G$
& $\ell$-进表示 $\mathrm{Gal}(\bar F/F)\to{}^{L}G$
& 平坦 ${}^{L}G$-丛 $E$（局部系统）
& 电特征膜 / 电磁对偶后的真空 \\
\midrule
自守侧对象
& 自守表示 of $G(\mathbb A_{\mathbb Q})$
& 自守表示 of $G(\mathbb A_F)$
& $D\text{-mod}(\mathrm{Bun}_G)$、Hecke 特征层 $\mathcal F_E$
& $B_{cc}$ 上的 $A$-膜范畴 \\
\midrule
对应
& Langlands 对应（猜想）
& Langlands 对应（函数域，$GL_n$ 已证）
& 几何 Langlands 对应
& 膜词典：$A\mapsto\operatorname{Hom}(B_{cc},A)$（argued） \\
\midrule
Hecke 算子
& Hecke 算子 on $G(\mathbb A)$
& Hecke 算子 on $G(\mathbb A_F)$
& Hecke 对应 $\mathrm{Bun}_G\xleftarrow{p}\mathcal H_z\xrightarrow{q}\mathrm{Bun}_G$
& 't Hooft 线算子 $T({}^{L}w,z_0)$ \\
\midrule
特征对象
& 与表示 $\pi$ 对应的 Hecke 特征函数
& 同左
& Hecke 特征层 $\mathcal F_E$：$H_R\mathcal F_E\cong V_R\otimes\mathcal F_E$
& 特征膜 eigenbrane（被 $W_R$ 作用成标量） \\
\midrule
几何舞台
& —（无）
& Jacobian / 曲线的模空间
& Hitchin 模空间 $\mathcal M_H(G,C)$，$\dim_{\mathbb C}=(2g-2)\dim G$
& 2d sigma 模型的目标空间 \\
\midrule
对偶机制
& 函数方程 / 互反律
& 函数方程
& 镜面对称 $A(\mathcal M_H({}^{L}G),\omega_K)\cong B(\mathcal M_H(G),J)$
& 电磁 $S$-对偶 + 紧化到 $C$ \\
\midrule
退化 / 特殊点
& —（无）
& —（无）
& Hitchin 纤维：复环面，复维 $(g-1)\dim G$
& $t=\pm i$ 复化、$t=\pm1$ 加乘、$t=0,\infty$ 退化 \\
\midrule
已验证的部分
& 类域论（$GL_1$）、$GL_2/\mathbb Q$ 局部情形
& $GL_n$（Drinfeld、Lafforgue）
& $GL_n$ 非分歧（Frenkel--Gaitsgory--Vilonen 等）
& 物理推导（KW06 §2--§11，argued） \\
\bottomrule
\end{longtable}
\endgroup

\vspace{4pt}
{\scriptsize\color{gcol}
读法：横看一行 = 同一件事；竖看一栏 = 一种语言。中列两栏是「已知 ↔ 未知」的桥：几何栏在 $GL_n$ 非分歧情形已是定理，
物理栏给出的是为什么该定理成立的机制（KW06 的贡献），但它本身是物理推导（argued）。
底本：arXiv:hep-th/0604151v3；对照：Witten 2008 (arXiv:0802.0999)、Frenkel 2009 (arXiv:0906.2747)。}

\end{document}
